\chapter*{Executive Summary}\label{ch:exec-sum}

\section{Motivation}
Modern biology has become increasingly data-rich but interpretation-limited. Data collection is now cost-effective in many realms, enabling the construction of novel datasets. High-throughput imaging technology has led to population-scale phenotyping datasets, while advanced sequencing methods allow researchers to capture genome-scale variation across populations. Global consortia efforts have standardized environmental metadata across climates and sites, while novel experimental data and computational frameworks have prompted the curation and creation of biological networks and large literature corpora. These data span fundamentally different structures: images contain spatial, pixel-level information; genotypes comprise high-dimensional, sparse signals; environments exist in a continuous and contextual space; and networks represent relational graph structures. Handling these different structures requires systems with different inductive biases; no single model class can handle them all.

\section{Thesis Statement}
While progress has been made applying deep learning to localized areas, biological data are heterogeneous and structurally varied, making the gains from current deep learning methods unequally distributed. Because of the nature of biological data, bridging the gap between inputs and interpretation requires problem-specific systems. Architectures, objectives, and benchmarks must be designed with appropriate inductive biases. This proposal demonstrates the utility of this framework across three linked stages of biological discovery: \textit{measurement}, \textit{prediction}, and \textit{interpretation.}

\section{Chapter 2 Summary}
Chapter 2 addresses the \textit{measurement} stage by extracting quantitative phenotypes from raw imagery. No deep learning framework currently exists for phenotyping pennycress (\textit{Thlaspi arvense L.}), an emerging oilseed cover crop. In this work, we present an automated deep learning pipeline that extracts 52 tissue-resolved phenotypes per seed pod. We benchmarked convolutional, transformer, and foundation-model segmentation architectures using 347 manually annotated pods from 21 scans, including an independent test set of 65 pods. A boundary-down-weighted U-Net achieved 85.58\% mean intersection over union while offering the most favorable accuracy--efficiency tradeoff. We then applied the pipeline to 12,253 seed pods from 768 accessions and connected the resulting trait matrix to candidate genetic mechanisms through heritability filtering, genome-wide association, network-based gene pruning, module construction, and expression analysis.

\section{Chapter 3 Summary}\label{sec:ch3}
Chapter 3 addresses the \textit{prediction} stage by mapping genotype and environment to phenotype. Accurate maize yield prediction is crucial for global food security, but remains challenging due to complex genotype-by-environment (G$\times$E) interactions. In this work, we present a unified transformer architecture for maize yield prediction that places 2,224 single-nucleotide variant markers and 702 environmental covariates in one sequence, allowing self-attention to learn direct G$\times$E interactions. A sparse mixture-of-experts layer supports token-level specialization, and an environment-conditioned affine head calibrates absolute yield while preserving within-environment rank. We trained on 143,050 usable pre-2024 Genomes2Fields plot records from 238 environments and retrospectively evaluated 9,486 records from 22 environments in the 2024 competition cohort. The selected model achieved a macro environment Pearson correlation coefficient of 0.423 and mean squared error of 55.40, corresponding to sixth place in the retrospective competition comparison and approaching the winning correlation of 0.45.

\section{Chapter 4 Summary}\label{sec:ch4}
Chapter 4 addresses the \textit{interpretation} stage by reasoning over biological networks to produce evidence-grounded mechanistic hypotheses. We propose \textsc{Mentor-RL}, an open-weight model trained on an ontology-grounded lattice of context-specific multiplex networks. A world-model curriculum teaches three co-equal structural views---random-walk neighborhoods, network hierarchies, and connectivity---before direct preference optimization and reinforcement learning train a tool-using policy with deterministic, schema-grounded rewards. Two flagship tasks test whether the model can generate a hierarchy for an unseen context and project a target neighborhood onto that hierarchy while recovering the network connections the hierarchy suppresses. Feasibility experiments demonstrate full fine-tuning of gpt-oss-120b on up to 128 Frontier nodes and an end-to-end trajectory loop. In preliminary S0 mapping experiments, atomic tokenization increased overall hard-subset accuracy from 45.6\% to 85.9\%, with both mapping directions reaching 100\% on the tested subset.

\section{Completion Plan and Expected Outputs}
Chapters 2 and 3 are being prepared for journal submission, while Chapter 4 defines the proposed remaining research and its staged evaluation plan. The Chapter 3 model and Chapter 4 feasibility experiments use distributed training on the OLCF Frontier supercomputer \cite{frontier}. Upon publication, we will release the open-source ground truth dataset of hand-annotated seed pod images and segmentation pipeline (Chapter~\ref{ch2:pennycress}); the training and evaluation code (Chapter~\ref{ch3:gxe}); and the training environment, evaluation harness, and trained checkpoints across pipeline stages (Chapter~\ref{ch4:mentor-rl}).
